\section{Covolume bounds}\label{sec:covolume}

\providecommand{\covolumelcv}{\operatorname{lcv}}
\providecommand{\covolumerk}{\operatorname{rk}}
\providecommand{\covolumemainT}{\mathrm{mainT}}
\providecommand{\covolumeunitE}{\mathrm{unitE}}
\providecommand{\covolumeII}{\mathrm{II}}
\providecommand{\covolumehdeg}{\widehat{\operatorname{deg}}}
\providecommand{\covolumehvol}{\widehat{\operatorname{vol}}}
\providecommand{\covolumeLat}[1]{\Lambda^{\mathrm{#1}}}

This section treats the proposition $\sfname{CovolumeBoundOfX}$ (\cref{def:covolX}), an archimedean
input of the hypothesis package of \cref{sec:families}, where it supplies the ``$L_3$ condition'' of the
indefinite-side data of the points. To each admissible point $p$ of a family and each index $j$ of the
chain of lifts there is attached a lattice $M_j\subseteq K^D$ over a totally complex field $K$, with
Hermitian metrics at the archimedean places, and $\sfname{CovolumeBoundOfX}$ asserts that eventually
\begin{equation}\label{eq:covolume:bound}
\tfrac1D\,\covolumelcv(M_j)\le\covolumemainT_j(p)+C\cdot\covolumeunitE(p),
\end{equation}
i.e.\ that the normalised Arakelov slope of $M_j$ is at least $-\covolumemainT_j(p)-C\,\covolumeunitE(p)$.
The main term is an integral multiple of the Chowla--Selberg-type constant
$\log\Omega=\gamma+\log(2\pi/\mathrm{rd}_E)-\Re\frac{L'}{L}(1,\chi_d)$ ($\mathrm{rd}_E$ the root discriminant
of the host field, $\chi_d$ the quadratic character modulo the auxiliary prime $d$ of \cref{sec:hosts}),
and $\covolumeunitE(p)$ is an error unit.

In the announcement \cite{Announcement2026} this part of the argument corresponds to the items ``Bost
slope inequality'' (item~13) and ``computation of covolumes'' (item~18), the Chowla--Selberg formula of
item~9 entering through $\log\Omega$. In the development the computation of covolumes is the bound
\eqref{eq:covolume:bound}, and the slope inequality (\cref{lem:covolume:slopes}) transports slopes along
maps of controlled height, from lattices of lifted forms to lattices of automorphic forms on unitary
Shimura varieties, where arithmetic intersection theory bounds them. The proof has four layers: a final
assembly (\cref{sec:covolume:assembly}), lift estimates (\cref{sec:covolume:lift}), mean slopes on deep
covers (\cref{sec:covolume:deep}) and arithmetic-geometric inputs (\cref{sec:covolume:inputs}), all of
them proved in the development.

% ---------------------------------------------------------------------------
\subsection{Hermitian lattices and slopes}\label{sec:covolume:lattices}

Let $K$ be a totally complex number field, $r_2=[K:\Q]/2$, and fix for each archimedean place $v$ an
embedding $\sigma_v\colon K\to\C$ inducing it.

\begin{definition}[Slopes and covolumes]\label{def:covolume:slopes}
Let $V$ be a $K$-vector space with a basis $(b_i)_{i\in I}$ and a \emph{frame} $X=(X(v))_v$ of
invertible complex $I\times I$ matrices. Put $\|x\|_{X,v}=|X(v)^{-1}(\sigma_v(x_i))_i|$ for
$x=\sum_ix_ib_i$, and let $\rho\colon V\to\bigoplus_v\C^I$ (Euclidean) be
$x\mapsto(X(v)^{-1}(\sigma_v(x_i))_i)_v$. For a nonzero finitely generated $\OO_K$-submodule
$L\subseteq V$ let $\covolumerk L=\dim_KKL$ and
\[
\mu(L)=-\frac{2\log\covol\rho(L)}{[K:\Q]\,\covolumerk L}-\log2+\frac{\log|d_K|}{[K:\Q]},\qquad
\deg L=\covolumerk L\cdot\mu(L).
\]
$\mu_{\max}(L)$ is the supremum of $\mu(N)$ over the nonzero submodules $N\subseteq L$, and
$\mu_{\min}(L)$ the infimum of $\mu(L/N)$ (quotient metrics) over the saturated $N\subsetneq L$. For
$V=K^d$ with its standard basis,
$\covolumelcv(X,L)=\frac2{[K:\Q]}\log\bigl(2^{r_2\covolumerk L}|d_K|^{-\covolumerk L/2}\covol\rho(L)\bigr)$.
The \emph{height} of a $K$-linear map $\varphi$ between such spaces is
$h(\varphi)=\frac1{r_2}\sum_v\log\|\varphi\|_v^2$ (operator norms). A family of norms
$\nu=(\nu_\sigma)_{\sigma\colon K\to\C}$ on $V$ \emph{has a Gram frame} if $\nu_{\sigma_v}=\|\cdot\|_{X,v}$
for all $v$, for some basis and frame; $\mu_\nu$, $\mu_{\max,\nu}$ are then defined through any such
frame. (The development sets $\mu_\nu=0$ if there is no Gram frame, and $\covol=0$ for non-discrete
groups; such cases are excluded by hypotheses wherever these quantities occur below.)
\end{definition}

\begin{remark}\label{rem:covolume:arakelov}
Let $\overline L$ be $L$ with the metrics $\|\cdot\|_{X,v}$ at $\sigma_v$ and $\bar\sigma_v$. Since
$\covol\rho(\OO_K)=2^{-r_2}|d_K|^{1/2}$ for the standard metric, a direct computation gives
$\mu(L)=2\covolumehdeg(\overline L)/([K:\Q]\covolumerk L)$ and, for $V=K^d$,
$\covolumelcv(X,L)=-\covolumerk L\cdot\mu(L)$, where $\covolumehdeg$ is the Arakelov degree
\cite{GilletSoule1992,Bost1996}. Thus $\mu$ is twice the normalised Arakelov slope, $\mu_{\max}$ and
$\mu_{\min}$ are the extreme slopes of the Harder--Narasimhan filtration \cite{Stuhler1976,Grayson1984},
and $\covolumelcv/\covolumerk$ is minus the slope.
\end{remark}

\begin{lemma}[Covolume versus slope]\label{lem:covolume:lcv-slope}
Let $\nu$ have a Gram frame, let $b_1,\dots,b_d\in V$ be linearly independent, $L\subseteq\sum_kKb_k$
an $\OO_K$-submodule, $X$ a frame on $K^d$ and $0<s\le1$ with
$\|c\|_{X,v}=s\,\nu_{\sigma_v}(\sum_kc_kb_k)$ for all $v$ and $c$. If
$M=\{c\in K^d:\sum_kc_kb_k\in L\}$ is a lattice, then $\covolumelcv(X,M)/d\le-\mu_\nu(L)$.
\end{lemma}

\begin{proof}
For $X'=sX$ the map $c\mapsto\sum_kc_kb_k$ is an isometry of $(M,X')$ onto $(L,\nu)$. Hence
$\covolumelcv(X',M)/d=-\mu_\nu(L)$ by \cref{rem:covolume:arakelov}, and passing from $X'$ to $X$
multiplies $\rho$ by $s$, whence $\covolumelcv(X,M)=\covolumelcv(X',M)+2d\log s$.
\end{proof}

The following facts are the form in which Bost's slope inequality \cite{Bost1996,Bost2001} and the
Hermite--Minkowski bounds of the geometry of numbers enter.

\begin{lemma}[Slope inequality; Minkowski bounds]\label{lem:covolume:slopes}
Let $A\subseteq V$ and $B\subseteq W$ be nonzero finitely generated $\OO_K$-submodules of spaces with
frames as in \cref{def:covolume:slopes}.
\begin{enumerate}[label=(\alph*)]
\item If $\varphi\colon V\to W$ is $K$-linear and injective and $\varphi(A)\subseteq B$, then\\
$\covolumerk A\,\bigl(\mu(A)-h(\varphi)\bigr)+(\covolumerk B-\covolumerk A)\,\mu_{\min}(B)\le\covolumerk B\cdot\mu(B)$.
\item If $B$ is spanned over $K$ by elements $y\in B$ with all $\|y\|_v\le e^\Lambda$, then
$\mu_{\min}(B)\ge-2\Lambda$.
\item If $\frac1{r_2}\sum_v\log\|y\|_v\ge-\Lambda$ for all $0\ne y\in B$, then
$\mu_{\max}(B)\le2\Lambda+\log\covolumerk B+\log|d_K|/[K:\Q]$.
\end{enumerate}
\end{lemma}

\begin{proof}
(a) If $B'$ is the saturation of $\varphi(A)$ in $B$, then $\deg B=\deg B'+\deg(B/B')$,
$\deg(B/B')\ge\covolumerk(B/B')\,\mu_{\min}(B)$, and
$\deg B'\ge\deg\varphi(A)\ge\deg A-\covolumerk A\cdot h(\varphi)$, as $\varphi$ multiplies covolumes at
$v$ by at most $\|\varphi\|_v^{2\covolumerk A}$. (b) One has $\mu(\OO_Ky)=-\frac1{r_2}\sum_v\log\|y\|_v^2$,
and every quotient of $B$ contains with full rank a module generated by images of the given
generators, whose degree is bounded below by Hadamard's inequality. (c) For $0\ne N\subseteq B$,
Minkowski's theorem gives $0\ne y\in N$ with $|\rho(y)|^2\le n\covol\rho(N)^{2/n}$,
$n=[K:\Q]\covolumerk N$; as $|\rho(y)|^2=\sum_v\|y\|_v^2$, the arithmetic-geometric mean inequality gives
$-2\Lambda\le\log(|\rho(y)|^2/r_2)\le\log\covolumerk N-\mu(N)+\log|d_K|/[K:\Q]$.
\end{proof}

% ---------------------------------------------------------------------------
\subsection{The statement}\label{sec:covolume:statement}

We use the families and points of \cref{sec:families}: for $\eps>0$ a family $F\in\mathrm{FamP}(\eps)$
provides, for $x>x_1$ and $\iota\in I_x$, global data (ranks $k,m,f,\dots$, a host field $E$ with
auxiliary primes $q\equiv3\bmod8$, $d\equiv7\bmod16$, a totally complex field $K$) and size data (the
scale $x$, the exponent $\eps$, parameters $b,e_z,\dots$), and $\mathrm{FamHyp}$ is the class of families
satisfying the family conditions. $\mathrm{Closed}(\mathcal C,P)$ means: there is $\eps_1>0$ such that for
every $\eps\in(0,\eps_1]$ there is a ``transported constant'' $\kappa$ with $P(F,C_H)$ for all
$F\in\mathrm{FamP}(\eps)\cap\mathcal C$ and $C_H\in\R$ such that the constants of $F$ and $C_H$ are large
relative to $\kappa$ ($F.\mathrm{Large}(\kappa,C_H)$). A \emph{point} $p\in\mathrm{Pt}(F,x,\iota)$ contains a
finite group $G_p$, a \emph{core datum} $T$ (unitary moduli spaces with lattices of automorphic forms),
spaces of lifted forms, and indefinite data $p.S(\mathrm{false})$, $p.S(\mathrm{true})$ over $K$ indexed by
$J=\{0,\dots,m-1\}$. $\mathrm{EventuallyAtAllPoints}(F,Q)$ means: for some $x_0$, $Q(p)$ holds for all
$x\ge x_0$, $\iota\in I_x$ and $p\in\mathrm{Pt}(F,x,\iota)$ satisfying the structural conditions
$\mathrm{PointConditions}'$ of \cref{sec:families} and $\mathrm{RecTied}$ (a bound for the degree and root
discriminant of an auxiliary number field, and a compatibility of two Teichm\"uller series).

\begin{definition}[Main term and error unit]\label{def:covolume:mainT}
For a point $p$ at $(x,\iota)$, with ranks $k,m,f$, host $E$, prime $d$ and parameters $b,\eps$, let
$\chi_d=(\frac{\cdot}{d})$, $\mathrm{rd}_E=|d_E|^{1/[E:\Q]}$, $v=\frac12\mathrm{rd}_E$, and
\begin{gather*}
\log\Omega(p)=\gamma+\log\frac{\pi}{v}-\Re\frac{L'(1,\chi_d)}{L(1,\chi_d)},\qquad
\covolumeunitE(p)=f\,m\,b\,x^{\eps}\log\log x,\\
\covolumemainT_j(p)=\sum_{\rho\in\mathrm{RPl}}\bigl((k+j)-(m-j-1)\bigr)\,w_{j+1}(\rho)\,\log\Omega(p)
\qquad(j\in J),
\end{gather*}
where $\mathrm{RPl}=\{0,1\}\times A$, with $A$ the finite abelian group of the datum (its elements label
elements of the CM type of the host), and $w_{j+1}(\rho)\in\Z$ is the entry $e_j$ of the
Harish-Chandra parameter of the datum at $\rho$.
\end{definition}

\begin{definition}[$\sfname{CovolumeBoundOfX}$]\label{def:covolX}
The base side $S=p.S(\mathrm{false})$ of a point contains an integer $D$ and, for $j\in J$, a frame
$X_{W,j}$ on $K^{d_W(j)}$ and an $\OO_K$-submodule $M_j\subseteq K^{d_W(j)}$. The proposition
$\sfname{CovolumeBoundOfX}$ is
\begin{multline*}
\mathrm{Closed}\Bigl(\mathrm{FamHyp},\,(F,C_H)\mapsto\exists C_0\ge0:\ \mathrm{EventuallyAtAllPoints}\bigl(F,\\
p\mapsto\forall C\ge C_0\ \forall j\in J:\ \tfrac1D\,\covolumelcv(X_{W,j},M_j)\le
\covolumemainT_j(p)+C\,\covolumeunitE(p)\bigr)\Bigr).
\end{multline*}
In the development the inner inequality is $S.\mathsf{L3}(p.N(C,C_H))$, where the numerical datum
$p.N(C,C_H)$ has $l_3(j)=\covolumemainT_j(p)+C\,\covolumeunitE(p)$.
\formal{GlobalTheory.CovTop.CovolumeBoundOfX}
\end{definition}

\begin{remark}\label{rem:covolume:omega}
(1) At the points considered $d_W(j)=D>0$ and $M_j$ is a lattice (\cref{lem:covolume:rec}), so the
inequality says $\mu_{X_{W,j}}(M_j)\ge-\covolumemainT_j(p)-C\,\covolumeunitE(p)$
(\cref{rem:covolume:arakelov}). $C_H$ enters only through $F.\mathrm{Large}(\kappa,C_H)$ and the choice of
$C_0,x_0$. (2) As $d\equiv3\bmod4$, $\chi_d$ is the odd character of $\Q(\sqrt{-d})$, and the functional
equation gives $\frac{L'}{L}(0,\chi_d)=\gamma+\log\frac{2\pi}{d}-\frac{L'}{L}(1,\chi_d)$; hence
$\log\Omega(p)=\frac{L'}{L}(0,\chi_d)+\log(d/\mathrm{rd}_E)$. By the Chowla--Selberg formula
\cite{ChowlaSelberg1967}, a consequence of Kronecker's limit formula, in the form given by Colmez
\cite{Colmez1993}, $\frac{L'}{L}(0,\chi_d)$ is an
explicit affine function of the Faltings height of the elliptic curves with complex multiplication by
the integers of $\Q(\sqrt{-d})$, i.e.\ of the logarithm of their period. The development uses
$\log\Omega$ only through its defining formula.
\end{remark}

% ---------------------------------------------------------------------------
\subsection{The main theorem}\label{sec:covolume:main}

\begin{theorem}[Covolume bound]\Fmark\label{thm:covol}
If the statements listed in \cref{tab:covolume:inputs} hold, then so does $\sfname{CovolumeBoundOfX}$.
Each of these fifteen statements is proved in the formal development; hence
$\sfname{CovolumeBoundOfX}$ holds.
\formal{SelfIntersectionVolumeEstimates.topCovolSlim2_use_alt_alt}
\end{theorem}

\begin{table}[htbp]
\centering\footnotesize
\begin{tabular}{@{}>{\raggedright\arraybackslash}p{2.8cm}>{\raggedright\arraybackslash}p{12.2cm}@{}}
\toprule
Input & Content (as formalised); literature\\
\midrule
Resolution & reduced quasi-compact schemes locally of finite type over a field of characteristic $0$
are resolved by blow-up sequences which are isomorphisms over the smooth locus, the reduced preimage of
the non-smooth locus having strict normal crossings \cite{Hironaka1964,BierstoneMilman1997,Wlodarczyk2005,Kollar2007}\\
Principalisation & an ideal sheaf on a smooth scheme becomes a divisor with strict normal crossings
support after blow-ups with centres transversal to the boundary (same references)\\
Poincar\'e--Lelong & arithmetic Poincar\'e--Lelong formula in multiprojective coordinates
\cite{GilletSoule1990,BostGilletSoule1994}\\
$\sfname{HSUse}$ & arithmetic Kleiman criterion, Hilbert--Samuel upper bounds (\cref{def:covolume:hs})
\cite{GilletSoule1992,AbbesBouche1995,Zhang1995}\\
Moriwaki & arithmetic self-intersection $\le$ arithmetic volume (\cref{def:covolume:hs})
\cite{Moriwaki2009,Yuan2008}\\
Hodge--Tate & exactness of the Hodge--Tate sequence of $p$-divisible groups over $\OO_C$
\cite{Faltings1988,FarguesGenestierLafforgue2008}\\
Zarhin & an abelian scheme over the valuation ring of an algebraically closed valued field is a
retract of a principally polarised one \cite{Zarhin1975,Zarhin1985}\\
Nef Riemann--Roch & asymptotic Riemann--Roch for nef divisors \cite{Lazarsfeld2004}\\
Siegel toroidal & toroidal compactifications of Siegel moduli spaces of odd level $n\ge3$ over
$\Z[\zeta_n,1/n]$, with semi-abelian extension, charts, ample Hodge bundle twist
\cite{FaltingsChai1990,AshMumfordRapoportTai1975}\\
Unitary toroidal & smooth toroidal models of unitary moduli spaces, charts, logarithmic
Kodaira--Spencer isomorphism \cite{Lan2013,FaltingsChai1990}\\
Lie seminegativity & subbundles of $\Lie G$, for $G$ semi-abelian over a projective curve in
characteristic $0$, have degree $\le0$ \cite{FaltingsChai1990,MoretBailly1985}\\
Point separation & point separation for adjoint bundles on klt pairs \cite{AngehrnSiu1995,Lazarsfeld2004}\\
$\sfname{ArDegree}$ & birational invariance of arithmetic self-intersections; Chern-form integrals
compute degrees \cite{GilletSoule1990}\\
Hodge package & $28$ statements (toroidal compactification, arithmetic volumes, vanishing theorems,
$L^2$ extension, Chow forms, heights of cycles, GAGA, \dots; \cref{sec:literature})\\
Grothendieck--Messing & deformations of abelian schemes over nilpotent divided-power thickenings
\cite{Messing1972}\\
\bottomrule
\end{tabular}
\caption{The inputs of \cref{thm:covol}; the first eight are those of \cref{thm:deepstr}.}
\label{tab:covolume:inputs}
\end{table}

\begin{proofoutline}
The first eight inputs give $\sfname{DeepStR}$ (\cref{thm:deepstr}). The others, with Green's theorem on
cuspidal characters (\cref{sec:literature}), a Fubini--Study/Chow height formula, a comparison of
arithmetic self-intersections with heights of Chow forms and the injectivity of small sections, are the
geometric inputs of \cref{thm:covolume:lift}, which gives $\sfname{CovolumeBoundOfX}$. In the final
proof the inputs are supplied by \cref{thm:covolume:resolution,thm:covolume:hs} and by formal proofs of
the other published results (\cref{sec:literature}). The part of the formal proof specific to
\cref{thm:covol} has about $850{,}000$ lines, more than half of them in the proof of \cref{thm:deepstr}.
\end{proofoutline}

% ---------------------------------------------------------------------------
\subsection{The final assembly}\label{sec:covolume:assembly}

Fix a point $p$ at $(x,\iota)$ and $i\in J$; write $\delta=\exp(-x^{e_z/7})$,
$\covolumemainT=\covolumemainT_i(p)$ and $U=\covolumeunitE(p)$. Let $W_i$ be the $i$-th space of lifted
forms of $p$, with the family $\nu$ of Petersson-type norms given by the archimedean realisations;
$\covolumeLat{all}\subseteq W_i$ the $\OO_K$-span of the integral structure;
$\covolumeLat{type}\subseteq\covolumeLat{all}$ the sublattice of vectors of the type prescribed by $p$;
$e_i$ the operator given by an idempotent of the group algebra of $G_p$;
$\covolumeLat{good}=\covolumeLat{type}\cap\operatorname{im}e_i$ and
$\covolumeLat{bad}=\covolumeLat{type}\cap\ker e_i$; and $\mu_{\mathrm{core}}$ the slope of the
corresponding lattice of automorphic forms of weight $\kappa_i$ of the core datum.

\begin{definition}[Lift conditions]\label{def:covolume:predicates}
The \emph{lift predicates} at $(p,i)$ are, for $A,B,C,PP\in\R$ and $|\kappa_i|$ the size of the weight
$\kappa_i$:
$\mathsf{GramAt}$: $\nu$ has a Gram frame, $\covolumeLat{all}$ is finitely generated and
$\covolumeLat{good}\ne0$ (and a condition on the weights of the core datum holds);
$\mathsf{RecAt}$: $\covolumelcv(X_{W,i},M_i)/D\le-\mu_\nu(\covolumeLat{good})$;
$\mathsf{CutAt}$: $e_i^2=e_i$, $e_i$ is \emph{Pythagorean}, i.e.\
$\nu_\sigma(e_iy+z-e_iz)^2=\nu_\sigma(e_iy)^2+\nu_\sigma(z-e_iz)^2$ for all $\sigma,y,z$,
$|G_p|\,e_i(\covolumeLat{type})\subseteq\covolumeLat{type}$, and $2\log|G_p|\le b$;
$\mathsf{BadAt}$: $\covolumerk\covolumeLat{bad}\le3\delta\,\covolumerk\covolumeLat{type}$;
$\mathsf{LamAt}(A)$: $\mu_{\max,\nu}(\covolumeLat{all})\le x^Ab$;
$\mathsf{SizeAt}(PP,B)$: $1\le b\le U\le x^Bb$, $|\covolumemainT|\le x^Bb$ and
$|\covolumemainT|,b\le PP\,|\kappa_i|$ (and conditions on auxiliary weights);
$\mathsf{CoreAt}(C)$: $\mu_{\mathrm{core}}-CU\le\mu_\nu(\covolumeLat{type})$;
$\mathsf{LowAt}(C)$: $-\covolumemainT-CU\le\mu_{\mathrm{core}}$.
The eight corresponding \emph{lift conditions} are the statements
$\mathrm{Closed}\bigl(\mathrm{FamHyp},(F,C_H)\mapsto\exists C\ge0:\mathrm{EventuallyAtAllPoints}(F,p\mapsto\forall i\in J:Q)\bigr)$,
the existential constant playing the role of $A$, $B$ or $C$, and $PP$ being a parameter of the family.
\end{definition}

\begin{lemma}[Good/bad cut]\label{lem:covolume:goodbad}
Let $V$ carry a frame, $e$ a $K$-linear idempotent of $V$ self-adjoint for the Hermitian forms at all
places, $g\ge1$ an integer, and $N\subseteq V$ a finitely generated $\OO_K$-submodule with
$g\,e(N)\subseteq N$. Put $N_g=N\cap\operatorname{im}e$, $N_b=N\cap\ker e$. Then
$\covolumerk N=\covolumerk N_g+\covolumerk N_b$, $\deg N\le\deg N_g+\deg N_b+2\covolumerk N_b\log g$, and
$\deg N_b\le\covolumerk N_b\cdot\mu_{\max}(M)$ for every finitely generated $M\supseteq N$.
\end{lemma}

\begin{proof}
For $y\in N$, $g\,e(y)\in N_g$ and $g(y-e(y))\in N_b$; hence $KN=KN_g\oplus KN_b$, which gives the
ranks, and $N_g=N\cap KN_g$ is saturated. As $e$ is self-adjoint, $KN_g\perp KN_b$, so the quotient norm
of $y+N_g$ is $\|(1-e)y\|$, and $g(1-e)$ induces an injection $N/N_g\hookrightarrow N_b$ of lattices of
equal rank, of height $2\log g$. Thus $\deg N=\deg N_g+\deg(N/N_g)$ and
$\deg(N/N_g)\le\deg N_b+2\covolumerk N_b\log g$ by \cref{lem:covolume:slopes}(a). The last claim is the
definition of $\mu_{\max}$.
\end{proof}

\begin{lemma}[Slope of the good part]\label{lem:covolume:slope-good}
Let $\nu$ have a Gram frame, $0\ne L\subseteq M$ with $M$ finitely generated, $e$ a Pythagorean
idempotent, $g\ge1$ an integer with $g\,e(L)\subseteq L$, and $0\le\delta\le\frac1{12}$ with
$\covolumerk(L\cap\ker e)\le3\delta\covolumerk L$. If $\mu_{\max,\nu}(M)\le Q$, $-Q\le\mu_\nu(L)$ and
$2\log g\le Q$, then $\mu_\nu(L)-12\delta Q\le\mu_\nu(L\cap\operatorname{im}e)$.
\end{lemma}

\begin{proof}
A Pythagorean idempotent is self-adjoint, so \cref{lem:covolume:goodbad} applies to $N=L$. Let
$d_g,d_b$ be the ranks of $L_g=L\cap\operatorname{im}e$, $L_b=L\cap\ker e$, $\mu=\mu_\nu(L)$,
$\mu_g=\mu_\nu(L_g)$ and $l=2\log g\in[0,Q]$. Then $(d_g+d_b)\mu\le d_g\mu_g+d_b(Q+l)$, and
$d_b\le3\delta(d_g+d_b)$ gives $d_b\le4\delta d_g$, so $d_g>0$. Using $\mu\ge-Q$,
$d_g\mu_g\ge(d_g+d_b)\mu-d_b(Q+l)\ge d_g\mu-3d_bQ\ge d_g(\mu-12\delta Q)$.
\end{proof}

\begin{theorem}[Final assembly]\label{thm:covolume:assembly}
The eight lift conditions imply $\sfname{CovolumeBoundOfX}$.
\formal{GlobalTheory.CovTop.Assembly.covolX_of_lift_bounds_and_coreLe}
\end{theorem}

\begin{proof}
$\mathrm{Closed}$ and $\mathrm{EventuallyAtAllPoints}$ are stable under finite conjunctions (take the
least $\eps_1$, a common transported constant and the largest $x_0$). Fix $\eps,\kappa,F,C_H$ as in
$\mathrm{Closed}$ and let $A$, $B$, $C_1\ge0$, $C_2\ge0$ be the constants of the lift conditions for
$\mathsf{LamAt}$, $\mathsf{SizeAt}$, $\mathsf{CoreAt}$ and $\mathsf{LowAt}$. Since $e_z>0$ is fixed and
$x^a\exp(-x^{e_z/7})\to0$, for large $x$ we have $12\delta Q\le U$, where
$Q=(x^A+x^B+C_1+C_2+1)U$. By $\mathsf{SizeAt}$, $1\le b\le U$ and $|\covolumemainT|\le x^BU$. Apply
\cref{lem:covolume:slope-good} to $L=\covolumeLat{type}\subseteq M=\covolumeLat{all}$, $e=e_i$ and
$g=|G_p|$: the Gram frame, finite generation and $L\ne0$ come from $\mathsf{GramAt}$; the conditions on
$e$ and $2\log g\le b\le Q$ from $\mathsf{CutAt}$; the rank condition is $\mathsf{BadAt}$;
$\mu_{\max}(M)\le x^Ab\le Q$ is $\mathsf{LamAt}$; and
$\mu(L)\ge\mu_{\mathrm{core}}-C_1U\ge-\covolumemainT-(C_1+C_2)U\ge-Q$ by $\mathsf{CoreAt}$ and
$\mathsf{LowAt}$. Hence $\mu(\covolumeLat{good})\ge\mu(L)-U$, and by $\mathsf{RecAt}$, $\mathsf{CoreAt}$,
$\mathsf{LowAt}$,
\[
\tfrac1D\covolumelcv(X_{W,i},M_i)\le-\mu(\covolumeLat{good})\le-\mu(L)+U\le-\mu_{\mathrm{core}}+(C_1+1)U
\le\covolumemainT+(C_1+C_2+1)U.
\]
So $C_0=C_1+C_2+1$ works, because $U\ge1$.

%P2PTR thm:covolume:assembly
A complete proof is given in \cref{sec:pf:lift2,sec:pf:covol}.
\end{proof}

% ---------------------------------------------------------------------------
\subsection{The lift estimates}\label{sec:covolume:lift}

The lift conditions are obtained by comparison with Shimura varieties of higher level. A \emph{deep
cover} $Y$ of the core datum at level $j$ is a unitary moduli space of higher level (level modulus $M'$)
over a coefficient field $c_2.F$ containing that of $T$, with complex uniformisations and dictionaries
to the local groups. It carries lattices $\mathrm{latW}(Y,G)$ of cusp forms of weight $G$, with
Petersson-type norms and slopes $Y.\mu(G)$, and, on each connected component $Z$, degrees, ranks,
arithmetic degrees and heights of automorphic bundles. $\mathrm{GoodCov}(Y)$ requires, among other
conditions, an admissible level, unitary uniformisations, normality of $c_2.F$ over the field of $T$, and
$\log[c_2.F:\Q]+\log|d_{c_2.F}|/[c_2.F:\Q]\le391z^2$ for a size parameter $z$. The weights used are
$G_k=G_k(\kappa_i)$, its shifts $G_s(\pm c_0)$ by a multiple of the Hodge weight, and an anchored weight
$u_M$.

\begin{theorem}[Lift estimates]\Fmark\label{thm:covolume:lift}
Assume $\sfname{DeepStR}$ (\cref{def:covolume:deepstr}) and the geometric inputs of the proof of
\cref{thm:covol}. Then the conditions of \cref{tab:covolume:conditions} hold; in particular the eight
lift conditions hold, and $\sfname{CovolumeBoundOfX}$ follows by \cref{thm:covolume:assembly}.
\formal{SelfIntersectionVolumeEstimates.topCovolG_alt_alt}
\end{theorem}
%P2PTR thm:covolume:lift
\begin{proof}See \cref{sec:pf:lift2}.
\end{proof}


\begin{table}[htbp]
\centering\footnotesize
\begin{tabular}{@{}>{\raggedright\arraybackslash}p{2.1cm}>{\raggedright\arraybackslash}p{9.1cm}>{\raggedright\arraybackslash}p{3.4cm}@{}}
\toprule
Group & Conditions at $(p,i)$, resp.\ at good deep covers $Y$ of level $i+1$ & Main inputs\\
\midrule
Point data & $\mathsf{GramAt}$; $\mathsf{SizeAt}$; $\nu_\sigma=c_\sigma\cdot$(Petersson norm) with
$|\log c_\sigma|\le CU$; $\covolumemainT$ equals an anchored sum $\sum_\tau(p-q)a(\tau)\log\Omega$ up to
$CU$; $\covolumeLat{type}$ spanned by integral vectors & point conditions, Hodge package\\
Deep covers & existence of $Y$, positive ranks, finitely many components, flat degrees; degrees and
rank shares of components; mean height of the components $\ge-X$ with $X\le CU$; Hilbert--Samuel lower
bounds; $2\sum_Z\covolumehdeg_Z/\sum_Z\covolumerk_Z\le Y.\mu(u_M)$; $Y.\mu(u_M)\ge-C_UU$ &
$\sfname{DeepStR}$, toroidal inputs, Grothendieck--Messing, Lie seminegativity, $\sfname{ArDegree}$\\
Descent & $Y.\mu(G_k)-2b\le\mu_{\mathrm{core}}$; jointly injective
$\varphi_h\colon\mathrm{latW}(Y,G_s(-c_0))\to\mathrm{latW}(Y,G_k)$, $\sum_h\|\varphi_hf\|^2\le
e^{2b}\|f\|^2$; a rank comparison; generators of norm $\le e^{PP|\kappa_i|}$;
$Y.\mu(u_M)-\covolumemainT-C_FU\le Y.\mu(G_s(-c_0))$ & Hodge package, rows above\\
Lift & $\mathsf{LowAt}$ (\cref{lem:covolume:low}; directly if $i+1=m$), $\mathsf{RecAt}$
(\cref{lem:covolume:rec}), $\mathsf{CutAt}$, $\mathsf{BadAt}$, $\mathsf{LamAt}$ (by
\cref{lem:covolume:slopes}(c)), $\mathsf{CoreAt}$ & rows above, Green\\
\bottomrule
\end{tabular}
\caption{The conditions established in the proof of \cref{thm:covolume:lift}, in logical order;
$U=\covolumeunitE(p)$, and $C,C_F,C_U$ are constants.}
\label{tab:covolume:conditions}
\end{table}

\begin{proofoutline}
Each row is a theorem of the development, proved from the indicated inputs and the rows above. The
bound $\sfname{DeepStR}$ enters only through the heights, and an arithmetic Hilbert--Samuel estimate on
each component turns the height, degree, share and sum bounds into $Y.\mu(u_M)\ge-C_UU$. Some step
theorems carry hypotheses that their proofs do not use; e.g.\ the Gram condition follows from the point
conditions alone. Three representative steps follow.
\end{proofoutline}

\begin{lemma}\label{lem:covolume:rec}
If $p$ satisfies $\mathrm{PointConditions}'$ and $\mathsf{GramAt}$ holds at $(p,i)$, then so does
$\mathsf{RecAt}$.
\end{lemma}

\begin{proof}
The primitive-rows condition gives $d_W(i)=D>0$ and that $M_i$ is a lattice. By construction of the
base side, $M_i=\{c\in K^D:\sum_kc_kb_k\in\covolumeLat{good}\}$ for a basis $(b_k)$ of the space of lifted
forms of the given type, and $\|c\|_{X_{W,i},v}=s\,\nu_{\sigma_v}(\sum_kc_kb_k)$ with $s=n_P^{-1/2}$ for an
integer $n_P\ge1$. Apply \cref{lem:covolume:lcv-slope}.
\end{proof}

\begin{lemma}[Descent]\label{lem:covolume:descent}
Let $A\subseteq V$, $B\subseteq W$ be finitely generated, $A\ne0$, with norm families having Gram
frames, and let $\varphi_1,\dots,\varphi_n\colon V\to W$ be jointly injective with
$\varphi_h(A)\subseteq B$ and $\sum_h\|\varphi_hf\|_v^2\le e^{2P}\|f\|_v^2$. Suppose that $B$ is spanned
over $K$ by elements of norm $\le e^\Lambda$ at all places, $\Lambda,\eps\ge0$,
$n\covolumerk B\le(1+\eps)\covolumerk A$ and $t\le\mu(A)-2P$. Then $\mu(B)\ge t-\eps(|t|+2\Lambda)$.
\end{lemma}

\begin{proof}
Apply \cref{lem:covolume:slopes}(a),(b) to $(\varphi_h)_h\colon A\to B^n$: its height is at most $2P$,
$\mu(B^n)=\mu(B)$ and $\mu_{\min}(B^n)\ge-2\Lambda$. With
$\theta=\covolumerk A/(n\covolumerk B)\in[(1+\eps)^{-1},1]$ this gives
$\mu(B)\ge\theta t-2\Lambda(1-\theta)\ge t-(1-\theta)(|t|+2\Lambda)$, and $1-\theta\le\eps$.
\end{proof}

\begin{lemma}\label{lem:covolume:low}
Let $C_F,C_U\ge0$ and $i+1<m$. If $\mathsf{SizeAt}(PP,\cdot)$ holds at $(p,i)$ and the descent
conditions and $Y.\mu(u_M)\ge-C_UU$ hold for a good deep cover $Y$ of level $i+1$, then
$\mathsf{LowAt}(2(C_F+C_U)+9)$ holds at $(p,i)$.
\end{lemma}

\begin{proof}
Let $\Lambda=PP|\kappa_i|\ge b\ge1$, $\eps=b/\Lambda$, $A=\mathrm{latW}(Y,G_s(-c_0))$,
$B=\mathrm{latW}(Y,G_k)$. Flatness and the uniform bound give
$t:=-\covolumemainT-(C_F+C_U)U-2b\le\mu(A)-2b$, and \cref{lem:covolume:descent} (with $P=b$) gives
$\mu(B)\ge t-\eps(|t|+2\Lambda)\ge-\covolumemainT-2(C_F+C_U)U-7b$, using $|\covolumemainT|\le\Lambda$
and $b\le\min(\Lambda,U)$. Finally $\mu_{\mathrm{core}}\ge\mu(B)-2b$ and $b\le U$.
\end{proof}

% ---------------------------------------------------------------------------
\subsection{Mean slopes on deep covers}\label{sec:covolume:deep}

\begin{definition}[$\sfname{DeepStR}$]\label{def:covolume:deepstr}
Let $X^o$ be a unitary moduli space attached to a CM field $E$ with CM type $\Psi$, a datum of signature
$(p,q)$, a level and a coefficient datum. For a weight $u=(u_{\tau,a})_{\tau\in\Psi,a<p+q}$ and $g>0$
let $\mathrm{WtGap}(u,g)$ mean that every $u_\tau$ is strictly decreasing with
$|u_{\tau,a}-u_{\tau,a'}|\ge g$ for $a\ne a'$, $a\ge p$; put
$\mathrm{bnd}(u,g)=4|\Psi|(p+q)^2(1+\max|u_{\tau,a}|/g)$ and
$\mathrm{err}_s(u,g)=(\log\mathrm{bnd}(u,g)+s)\|u\|_1+q\sum_\tau(\max_au_{\tau,a}-\min_au_{\tau,a})$.
$\mathrm{MeanAll}(s,C)$ means: whenever $\mathrm{WtGap}(u,g)$, there are reals
$\mathrm{sl}(Z)\le\mathrm{volSlope}(Z,u)$ for the connected components $Z$ with
$\frac1{\#\{Z\}}\sum_Z\mathrm{sl}(Z)\ge-C\,\mathrm{err}_s(u,g)$, where $\mathrm{volSlope}(Z,u)$ is the
arithmetic volume slope on $Z$ of the automorphic line bundle of weight $u$ (counting sections small at
the archimedean places, with valuation conditions at the primes up to $\mathrm{bnd}(u,g)$).
$\sfname{DeepStR}$ asserts: there is $C\ge0$ such that for every global datum $\mathcal D$, every core
datum $T$ of $\mathcal D$ satisfying the structural conditions $\mathsf{LawsJ}$, with $m^2\le k$, every $1\le j<m$ and every
deep cover $Y$ of $T$ at level $j$ that is good for some parameters and whose level modulus has at most
$15$ prime factors, $\mathrm{MeanAll}(\ell(\mathcal D),C)$ holds on $Y$, where
$\ell(\mathcal D)=\log\mathrm{rd}_E+|\frac{L'}{L}(1,\chi_d)|$.
\formal{GlobalTheory.CompanionTheory.Ar.Plug.DeepStR}
\end{definition}

\begin{theorem}[Mean slopes on deep covers]\Fmark\label{thm:deepstr}
Resolution and principalisation in characteristic $0$, the arithmetic Poincar\'e--Lelong formula,
$\sfname{HSUse}$, $\sfname{MoriwakiBigStatement}$, the exactness of the Hodge--Tate sequence, Zarhin's
trick and asymptotic Riemann--Roch for nef divisors imply $\sfname{DeepStR}$.
\formal{PositivityTheory.HtPt.Ah.deep_slim_use}
\end{theorem}

\begin{proofoutline}
(1)~Resolution and principalisation give resolution in projective coordinates, and with the
Poincar\'e--Lelong formula a full arithmetic intersection system (\cref{def:covolume:hs}); the other
inputs of \cref{thm:tower} (homomorphisms of semi-abelian schemes, GAGA for sections, quotients by
finite flat group schemes) are theorems of the development. (2)~One derives lines of large degree and
fourteen geometric statements (among them the birational pull-back and push-forward inequality for
arithmetic volumes), and then \cref{thm:tower}. (3)~The mean-slope part of $\sfname{VolumeTheoremX}$
gives $C$ such that $\mathrm{MeanAll}(\log\mathrm{rd}_E+|\frac{L'}{L}(1,\chi_d)|,C)$ holds for every abelian CM
field $E$, datum with $q\ge1$ and $3q+1\le p+q$, level with good reduction everywhere and trivial
$A$-part, and at most $4|\Psi|(p+q)^2$ bad primes. A good deep cover at level $j$ has $q=m-j\ge1$,
$p+q=k+m-1$ and $|\Psi|=f\ge2$, so $3q+1\le p+q$ follows from $m^2\le k$; its bad primes are controlled
by the prime factors of the level modulus.

%P2PTR thm:deepstr
A complete proof is given in \cref{sec:pf:deepstr}.
\end{proofoutline}

\begin{definition}[Root towers; the tower statements]\label{def:covolume:tower}
A \emph{root tower} of type $(\Psi,N)$ consists of $1\le j<N$, finitely many components $c$, a finite set
$\mathrm{bad}\ni2$ of primes and, for $B\in\N$ and each $c$, a model realising flag bundles over $c$, with
a geometric degree functional and, for every intersection system, arithmetic intersection functionals
$\mathrm{PP}(B,c,\cdot)$ on polynomials in the Chern roots; one forms volumes $\vol(c)$ and slopes
$\mathrm{slope}(B,c,u)=\mathrm{PP}(B,c,\mathrm{lin}(u)^{n+1})/((n+1)\vol(c)\mathrm{flagDeg}(u))$, $n$ the
dimension of the flag variety. $\sfname{GoodTowerExistsX}$: for every CM field, datum with $p,q\ge1$,
level and global model with split components there is a root tower of type $(\Psi,p+q)$ with a
\emph{good origin} (smooth flat models of flag bundles, ample reference models) satisfying a local
holomorphy condition. $\sfname{VolumeTheoremX}$: there is $C\ge0$ such that for every \emph{abelian} CM
field, every such tower and every full intersection system, if $q\ge1$, $3q+1\le p+q$ and
$\#\mathrm{bad}\le4|\Psi|(p+q)^2$, then: the geometric degree functional is $\vol(c)$ times the flag
degree, with $\vol(c)>0$ independent of $c$ and $B$ (proportionality); limit formulas (as the shift of
the bundles tends to $0$) and an asymptotic formula (up to $O(\log B/B)$) hold for the arithmetic
intersection functionals; and for weights $u$ with gap $g$, $B\ge128$ large in terms of $(u,g)$ and
$\#\mathrm{bad}\le B$,
\[
\frac1{\#\{c\}}\sum_c\mathrm{slope}(B,c,u)\ge-C\Bigl(\bigl(\log(B\,\mathrm{rd}_E)+\bigl|\tfrac{L'}{L}(1,\chi_d)\bigr|\bigr)\|u\|_1+\min(j,N-j)\,\mathrm{spread}(u)\Bigr),
\]
with a corresponding bound for $\sum_c\mathrm{PP}(B,c,\mathrm{lin}(u)^{n+1})$.
\end{definition}

\begin{theorem}[Good tower volume theorem]\Fmark\label{thm:tower}
The existence of a full arithmetic intersection system, $\sfname{HSUse}$, $\sfname{MoriwakiBigStatement}$,
resolution in projective coordinates, the extension of homomorphisms of semi-abelian schemes, the
Hodge--Tate statement, GAGA for sections, the existence of lines of large degree, quotients by finite
flat group schemes and Zarhin's trick imply $\sfname{GoodTowerExistsX}\wedge\sfname{VolumeTheoremX}$.
\formal{GlobalTheory.CompanionTheory.Ar.Link.good_tower_volume_theorem_bis}
\end{theorem}

\begin{proofoutline}
Further inputs are the integrability of points (from resolution), an avoidance statement, an input from
the Weil conjectures, and a positivity statement for families of abelian schemes (from the Hodge--Tate,
quotient and Zarhin inputs). The existence part constructs root towers with good origins; the volume
part combines an antisymmetry identity for good towers, bounds for test heights, a regularisation of
cusp weights, the equality of volumes given by a full intersection system, limits of mixed integrals off
divisors, and nef and height bounds for the Hodge bundle and the Lie metric. Proportionality is
Hirzebruch--Mumford proportionality in the non-compact case \cite{Mumford1977}; we know of no published
source for $\sfname{VolumeTheoremX}$ or $\sfname{DeepStR}$ in this generality (uniformly over abelian CM
fields, levels and signatures with $p\ge2q+1$), and they are new results of the development.

%P2PTR thm:tower
A complete proof is given in \cref{sec:pf:tower}.
\end{proofoutline}

% ---------------------------------------------------------------------------
\subsection{Arithmetic-geometric inputs}\label{sec:covolume:inputs}

\begin{definition}[Arithmetic varieties in coordinates]\label{def:covolume:hs}
Let $R=\Z[x_{k,i}:k\in\mathcal K,\,i\in\iota_k]$ with its $\Z^{\mathcal K}$-grading ($\mathcal K$ and the
$\iota_k$ finite).
(a)~An \emph{arithmetic variety} of dimension $d$ is a relevant multihomogeneous prime $P\subset R$ with
$P\cap\Z=0$ and $\dim R/P=d+|\mathcal K|$, i.e.\ the multicone of an integral closed subscheme of
$\prod_k\mathbb P^{\iota_k}_\Z$, flat over $\Z$, of absolute dimension $d$.
(b)~A \emph{Hermitian line} $A$ is a multidegree $a$ with a positive function $h$ on the complex
multicone, $h(c\cdot z)=\prod_k|c_k|^{a_k}h(z)$; a section $s$ of degree $a$ has $\|s\|=|s|/h$;
smoothness, semipositivity and nefness (on fibres, on the generic fibre) are defined in these terms.
The \emph{arithmetic volume} $\covolumehvol(A)$ is the $\limsup$ of
$\log\#\{\text{small sections of degree }\tau a\}/(\tau^{D+1}/(D+1)!)$, small sections being elements of
the function field integral on the standard charts with supremum norm $\le1$.
(c)~An \emph{intersection family} assigns numbers $\covolumeII(P;A_1,\dots,A_n)$, symmetric and additive,
with the recursion
$\covolumeII(P;A_1,\dots,A_n,A)=\sum_W\mathrm{mult}_W(P+(s))\,c_W-\int_{P(\C)}\log\|s\|\,c_1(A_1)\cdots c_1(A_n)$
for sections $s\notin P$ of the degree of $A$, the sum over the components $W$ of $P+(s)$, where
$c_W=\covolumeII(W;A_1,\dots,A_n)$ if $W$ is horizontal and $c_W$ is the geometric degree times $\log\ell$
if $W$ lies over the prime $\ell$; it is \emph{full} if it is moreover multilinear and satisfies the
projection formula and the expected vanishing.
(d)~$\sfname{HSUse}$ is the conjunction of the \emph{arithmetic Kleiman criterion} (if $\covolumeII(W;A)\ge0$
for all arithmetic curves $W\supseteq P$, then $\covolumeII(P;A,\dots,A)\ge0$) and two \emph{Hilbert--Samuel
upper bounds} $\covolumehvol(A)\le\covolumeII(P;A^{D+1})$ ($P$ of dimension $D+1$, $A$ conjugation
invariant, semipositive and nef on the fibres, plus technical conditions on the section lattices).
$\sfname{MoriwakiBigStatement}$ is
$\covolumeII(P;A^{D+1})\le\covolumehvol(A)$ for such $A$ which are moreover nef on the generic fibre.
Both quantify over all intersection families.
\end{definition}

The recursion in (c) is the inductive formula of Bost, Gillet and Soul\'e for heights
\cite{BostGilletSoule1994}; it determines $\covolumeII$, so an intersection family is the arithmetic
intersection theory of \cite{GilletSoule1990,BostGilletSoule1994}, and its existence amounts to the
arithmetic Poincar\'e--Lelong formula (independence of $s$). Besides the arithmetic Kleiman criterion
\cite{Zhang1995,Moriwaki2009}, (d) contains the upper-bound half of the arithmetic Hilbert--Samuel formula
\cite{GilletSoule1992,AbbesBouche1995,Zhang1995} and Moriwaki's inequality
$\covolumehvol(\overline L)\ge\covolumehdeg(\overline L^{D+1})$ for $\overline L$ nef on the fibres and on the
generic fibre \cite{Moriwaki2009,Yuan2008}. The formal statements are variants of the published ones
(explicit coordinates, sections counted in the function field, metrics given by homogeneous functions,
an arbitrary intersection family); for smooth metrics they contain the classical statements used here.

\begin{theorem}[Hilbert--Samuel and Moriwaki in coordinates]\Fmark\label{thm:covolume:hs}
Resolution of singularities in projective coordinates, the existence of a full intersection system and
the Zariski density of arithmetic curves of bounded height imply the statements of
\cref{def:covolume:hs}(d). These three hypotheses are theorems of the development (by
\cref{thm:covolume:resolution}, the formal arithmetic Poincar\'e--Lelong formula, and a height comparison
for orbit ideals); hence $\sfname{HSUse}$ and $\sfname{MoriwakiBigStatement}$ hold.
\formal{PositivityTheory.Pf1.HilbertSamuelMoriwaki.hsChain}
\end{theorem}

\begin{proofoutline}
A chain of $180$ steps: Kleiman's ampleness criterion \cite{Kleiman1966}; nefness and bigness on fibres;
counting small sections along multiples, lifting sections to resolved models; the arithmetic volume of
big lines via separation of sections by points of bounded height \cite{Moriwaki2009,Yuan2008}; curves by
lattice-point counting; arithmetic Kleiman by integration against positive forms; and the
Hilbert--Samuel upper bound via successive minima and $L^2$-lattices of sections
\cite{GilletSoule1992,AbbesBouche1995,Zhang1995}.
\end{proofoutline}

\begin{theorem}[Resolution and principalisation {\cite{Hironaka1964,Villamayor1989,BierstoneMilman1997,Wlodarczyk2005,Kollar2007}}]\Fmark\label{thm:covolume:resolution}
Let $k$ be a field of characteristic $0$. (a)~There is a rule assigning to every reduced quasi-compact
scheme $X$ locally of finite type over $k$ a blow-up sequence $X'\to X$, proper and an isomorphism over
the smooth locus, with $X'$ smooth over $k$ and centres over the non-smooth locus, such that the
reduced preimage of the non-smooth locus has strict normal crossings. (b)~There is a rule assigning to
every smooth quasi-compact $X$ over $k$, ideal sheaf $I$ nonzero on every component and boundary with
simple normal crossings a blow-up sequence with smooth centres, transversal to the accumulated boundary
and inside the support of the transformed ideal, after which $I\OO_{X'}$ is a product of exceptional
ideals. Both rules commute with smooth morphisms of constant relative dimension and with extensions of
$k$. In particular the first two inputs of \cref{tab:covolume:inputs} hold.
\formal{FunctorialResolution.functorialResolutionCharZero}
\formal{FunctorialResolution.functorialPrincipalizationCharZero}
\end{theorem}

\begin{proofoutline}
The development formalises marked ideals with boundary and \emph{principalisation rules}: order
reduction (coefficient ideal, maximal-contact hypersurfaces, descent), maximal-order reduction (a
companion-ideal reduction and principalisation of a monomial ideal) and divisor reduction (restriction
to a boundary divisor), each defined from the rules in lower dimension. The central theorem,
\texttt{canonicalAlgorithmRules\_\allowbreak isCorrect}, states that the canonical rules reduce inductively on
dimension: if the rules in lower dimension have the required properties (admissible centres, resolution
of the marked ideal, compatibility with field extensions and smooth pull-backs), so do the three
operations. Induction gives functorial principalisation, and gluing affine embedded principalisations
gives (a). This agrees closely with \cite{BierstoneMilman1997,Wlodarczyk2005,Kollar2007}; normal crossings
are strict, and smoothness of centres is not part of (a).
\end{proofoutline}

\emph{Toroidal compactifications and abelian schemes.} The inputs on Shimura varieties are formal
versions of the theory of Faltings and Chai \cite{FaltingsChai1990} and its extension to PEL-type, here
unitary, moduli spaces \cite{Lan2013}, built on \cite{AshMumfordRapoportTai1975}. Zarhin's trick
\cite{Zarhin1975,Zarhin1985} ($(A\times A^\vee)^4$ is principally polarisable), the Hodge--Tate sequence
\cite{Faltings1988,FarguesGenestierLafforgue2008} and quotients by finite flat group schemes enter
\cref{thm:tower} through the positivity of families of abelian schemes; the semi-negativity of Lie
algebras of semi-abelian schemes \cite{FaltingsChai1990,MoretBailly1985} and Grothendieck--Messing
theory \cite{Messing1972} enter the lift estimates. All are proved in the development
(\cref{sec:literature}).
